\documentclass[11pt,a4paper]{article}

% =============================================================================
% Pre-deployment compositional performance forecasting:
% problem statement, literature review, mathematical definition, and results.
% =============================================================================
\usepackage[utf8]{inputenc}
\usepackage[T1]{fontenc}
\usepackage[margin=0.88in]{geometry}
\usepackage{amsmath,amssymb,amsfonts,mathtools,bm}
\usepackage{booktabs,longtable,array,multirow}
\usepackage{graphicx}
\usepackage{xcolor}
\usepackage{enumitem}
\usepackage{microtype}
\usepackage{pdflscape}
\usepackage[colorlinks=true,linkcolor=blue!55!black,citecolor=green!40!black,urlcolor=blue!60!black]{hyperref}

\setlength{\parindent}{0pt}
\setlength{\parskip}{0.45em}
\setlist{leftmargin=1.5em,itemsep=2pt,topsep=3pt}
\allowdisplaybreaks
\newcolumntype{P}[1]{>{\raggedright\arraybackslash}p{#1}}

\newcommand{\R}{\mathbb{R}}
\newcommand{\E}{\mathbb{E}}
\newcommand{\acc}{\operatorname{acc}}
\newcommand{\MAE}{\operatorname{MAE}}
\newcommand{\clip}{\operatorname{clip}}
\newcommand{\sign}{\operatorname{sign}}
\newcommand{\softmax}{\operatorname{softmax}}
\newcommand{\pp}{\,\text{percentage points}}
\newcommand{\Rzero}{\mathrm{R0}}
\newcommand{\Rone}{\mathrm{R1}}
\newcommand{\Rtwo}{\mathrm{R2}}
\newcommand{\Rthree}{\mathrm{R3}}
\newcommand{\Rfour}{\mathrm{R4}}

\title{\textbf{Pre-Deployment Compositional Performance Forecasting}\\
\large Problem Definition, Mathematical Approach, Literature Review, and CIFAR Results}
\author{Technical research summary for the ShiftAlgebra repository}
\date{October 2026}

\begin{document}
\maketitle

\begin{abstract}
A deployed classifier can encounter several shifts at once: blur followed by
noise, low contrast followed by rotation, or a natural domain shift followed by
a sensor corruption.  Evaluating every ordered combination is combinatorially
expensive.  This project studies a specific earlier-stage forecasting problem:
predict the behavior of a frozen model under an unavailable composition using
measurements of primitive shifts and a limited number of ordered pairs, without
processing images from the composed target at forecast time.  This document
states that problem in plain and mathematical language, distinguishes it from
robust training and ordinary out-of-distribution (OOD) performance prediction,
reviews the closest literature, explains the matrix and hazard formulations,
and reports the complete empirical progression from the original CIFAR
experiments through the sealed Phase-2 model-bank pilot.  The record is mixed
and is reported as such: historical development experiments were positive,
independent Phase-1 and Phase-1B long-word tests rejected the original primary
method, while the Phase-2 pilot found that composition-aware hazard forecasts
substantially improve prediction and model selection across a bank of CIFAR
models.  The evidence supports a narrow composition-aware forecasting claim,
not a universal algebra of distribution shift.
\end{abstract}

\tableofcontents
\newpage

% =============================================================================
\section{The problem in plain language}
% =============================================================================

Suppose a cloud or vision system has already been trained.  Engineers can test
one fault or one shift at a time: blur, noise, low contrast, packet delay,
resource pressure, or sensor failure.  Real deployments can encounter several
of them together and in a meaningful order.  Testing every possible sequence
is expensive because the number of ordered words grows rapidly.  With $K$
primitive shifts, there are $2^K-1$ nonempty unordered subsets, and the number
of distinct ordered words of lengths $1$ through $K$ is
\[
\sum_{m=1}^{K}\frac{K!}{(K-m)!},
\]
which is even larger.

The practical question is:

\begin{quote}
\textbf{Can we forecast how an already-trained model will behave under a new
ordered combination of shifts before samples from that combination are
available?}
\end{quote}

The proposed workflow measures each simple shift and a controlled budget of
pairs.  It then learns a rule for combining those measurements.  At deployment
planning time, the rule receives the name/order of a new composition and the
stored primitive/pair profiles; it does \emph{not} run the unavailable target
composition through the classifier.

This is useful for pre-deployment risk analysis, test prioritization, and model
selection.  It does not make the classifier robust by itself.  It forecasts
where the classifier is likely to fail.

% =============================================================================
\section{Formal problem definition}
% =============================================================================

\subsection{Frozen model, primitive shifts, and words}

Let
\[
f:\mathcal{X}\rightarrow\Delta^{C-1}
\]
be a frozen $C$-class classifier, where $f(x)$ is a probability vector.  Let
$D_0$ be a labeled clean reference distribution.  A primitive shift $X$ is an
operator
\[
\tau_X:\mathcal{X}\rightarrow\mathcal{X}.
\]
Examples include contrast change, blur, noise, inversion, rotation, and swirl.
An ordered word
\[
w=(X_1,X_2,\ldots,X_m)
\]
means that $X_1$ is applied first:
\[
\tau_w(x)=\tau_{X_m}\!\left(\cdots
\tau_{X_2}(\tau_{X_1}(x))\cdots\right).
\]
Order can matter.  In general,
\[
\tau_{X\to Y}\neq\tau_{Y\to X},
\qquad
\acc_f(X\to Y)\neq\acc_f(Y\to X).
\]
For example, noise added after blur remains sharp, while blur applied after
noise can remove some high-frequency noise.

\subsection{Behavior to be forecast}

The simplest target is accuracy
\[
a_f(w)=\Pr_{(x,y)\sim D_0}
\left[\arg\max_c f_c(\tau_w(x))=y\right].
\]
The broader project also forecasts confidence margin, negative log likelihood,
expected calibration error, and representation drift.  Write the behavior
vector as
\[
B_f(w)=[\text{margin},\text{drift},\text{ECE},\text{NLL},\text{accuracy}],
\]
and the change from clean behavior as
\[
\Delta B_f(w)=B_f(w)-B_f(\varnothing).
\]

For every primitive $X$, the system can store:
\begin{itemize}
  \item a measured behavior vector $\Delta B_f(X)$;
  \item an architecture-independent fingerprint $v_f(X)$ containing fixed-size
  summaries of pixel, feature, logit, probability, displacement, and spectral
  changes.
\end{itemize}
For selected ordered pairs, it can additionally store
$\Delta B_f(X\to Y)$.

\subsection{Forecasting objective}

The primary target-free forecaster is
\begin{equation}
\widehat B_f(w)=
\mathcal{F}_{\Theta}\!\left(
  f\text{-metadata},
  \{v_f(X_i),\Delta B_f(X_i)\}_{i=1}^{m},
  \{\Delta B_f(X_i\to X_j)\}_{(i,j)\in\mathcal{C}},
  w
\right),
\label{eq:forecast}
\end{equation}
where $\mathcal{C}$ is a declared pair-calibration budget.  The target images
$\{\tau_w(x)\}$ and target labels are unavailable to
$\mathcal{F}_{\Theta}$ at forecast time.

The primary accuracy error is reported in percentage points:
\begin{equation}
\MAE_{\mathrm{acc}}=100\cdot
\frac{1}{|\mathcal{W}_{\mathrm{test}}|}
\sum_{w\in\mathcal{W}_{\mathrm{test}}}
\left|\widehat a_f(w)-a_f(w)\right|.
\end{equation}
When many classifiers are available, the task also asks whether
$\widehat a_f(w)$ correctly ranks models for a fixed environment $w$ and
selects a good deployment model.

\subsection{What ``zero-shot'' means here}

``Zero-shot'' refers only to the held-out composition.  It does \emph{not}
mean that the system uses no labels at all.  Primitive behaviors and selected
pair behaviors can be measured with labeled calibration data.  The claim is:
\begin{quote}
No images or labels from the held-out composed target enter the ordinary
\Rtwo{} forecast.
\end{quote}
This access boundary must be stated whenever results are compared with methods
that see actual unlabeled target images.

% =============================================================================
\section{Information regimes: fair comparisons require equal access}
% =============================================================================

\begin{center}
\begin{tabular}{P{1.2cm}P{5.4cm}P{4.6cm}P{3.5cm}}
\toprule
\textbf{Regime} & \textbf{Available at forecast time} & \textbf{Examples} & \textbf{Interpretation} \\
\midrule
\Rzero{} & Clean/source labels and model metadata; no primitive or target images & Clean accuracy, architecture/recipe ridge & Minimum source baseline \\
\Rone{} & \Rzero{} plus primitive profiles and primitive behavior & Mean, max damage, product retention, signed $L_p$ & Primitive-only forecast \\
\Rtwo{} & \Rone{} plus a declared set of labeled ordered pairs; no composed target images & Shift Algebra, pair residual, hazard composition, sequence baselines & Primary pre-deployment regime \\
\Rthree{} & Actual unlabeled composed target images are available & ATC, DoC, COT/COTT, MaNo, MDE, Agreement & Standard target-available OOD prediction \\
\Rfour{} & Labeled composed target is available & Direct measurement/oracle & Evaluation ceiling only \\
\bottomrule
\end{tabular}
\end{center}

The primary scientific comparison must be within \Rtwo{}.  An \Rthree{} method
can be operationally valuable, but it has information that an \Rtwo{} method
does not.  Therefore a claim such as ``better than ATC'' is meaningful only if
both information regimes are clearly labeled and the official implementation
is used.

% =============================================================================
\section{Two mathematical approaches used in this project}
% =============================================================================

\subsection{Shift Algebra: transformed behavior, symmetric interaction, order}

For behavior coordinate $b$ and exponent $p\geq 1$, define the exact signed
power transform
\begin{align}
\phi_p(b)&=\sign(b)\left((|b|+\epsilon)^p-\epsilon^p\right),\\
\phi_p^{-1}(u)&=\sign(u)
\left((|u|+\epsilon^p)^{1/p}-\epsilon\right).
\end{align}
The base composition of primitive behaviors is
\begin{equation}
B_{L_p}(w)=
\phi_p^{-1}\!\left(\sum_{i=1}^{m}\phi_p(\Delta B_f(X_i))\right).
\label{eq:lp}
\end{equation}
For effects with the same sign, $p>1$ lies between a maximum and an additive
sum, providing a simple saturation law.

Each primitive also receives a symmetric coordinate $z_i$ and directional
coordinates $(\pi_i,\kappa_i)$.  For a word of length $m$, the closed-form
features are
\begin{align}
\Gamma(w)&=\frac{1}{\binom{m}{2}}
\sum_{i<j}z_i\odot z_j,\\
\Omega(w)&=\frac{1}{\binom{m}{2}}
\sum_{i<j}\frac12\pi_i\odot\pi_j\odot
(\kappa_j-\kappa_i).
\label{eq:order}
\end{align}
The symmetric channel is unchanged by reversal, while the ordered channel
changes sign.  A decoder maps
$[B_{L_p}(w),\Gamma(w),\Omega(w)]$ to predicted behavior.

\subsection{Where the matrices come from}

The matrix structure is fixed analytically; training learns the coordinates
inserted into it.  For one directional coordinate, a primitive has the
Heisenberg block
\begin{equation}
H(u,v)=
\begin{pmatrix}
1 & u & \tfrac12uv\\
0 & 1 & v\\
0 & 0 & 1
\end{pmatrix},
\qquad u=\kappa\pi,\quad v=\pi.
\end{equation}
If Blur has $(\pi_B,\kappa_B)=(0.6,-0.3)$ and Noise has
$(\pi_N,\kappa_N)=(0.7,0.4)$, then
\[
H_B=
\begin{pmatrix}1&-0.18&-0.054\\0&1&0.6\\0&0&1\end{pmatrix},
\qquad
H_N=
\begin{pmatrix}1&0.28&0.098\\0&1&0.7\\0&0&1\end{pmatrix}.
\]
Chronological Blur then Noise is $H_NH_B$; reversing the word gives $H_BH_N$.
After removing the order-independent central part, the directional features are
$+0.147$ and $-0.147$.  Pair-supervision backpropagation changes the encoder,
$(\pi,\kappa)$, and decoder weights so that these features reduce measured pair
prediction error.

This is supervised behavior prediction, not an autoencoder.  There is no
objective that reconstructs the original image.  Also, the matrix extractor is
exactly equivalent to the closed-form all-pairs expressions in
Equations~\eqref{eq:lp}--\eqref{eq:order}; the matrix form supplies an
associative streaming state, not a different predictor by itself.

\subsection{Why the original algebra failed on long calibrated words}

Phase 1B showed that averaging pair interactions by $\binom{m}{2}$ can make the
correction too weak when damage accumulates extensively with length.  The
method can be stable but systematically optimistic.  That negative result led
to the Phase-2 hazard formulation.

\subsection{Phase-2 extensive hazard formulation}

Let $a_0$ be clean accuracy and $c=1/C$ chance accuracy.  Define normalized
retention and hazard
\begin{align}
r(a)&=\clip\!\left(\frac{a-c}{a_0-c},0.02,2\right),\\
h(a)&=-\log r(a).
\end{align}
For primitive $i$, write $h_i=h(a_i)$.  For ordered pair $i\to j$, define the
pair residual
\begin{equation}
R_{ij}=h(a_{i\to j})-h_i-h_j.
\end{equation}
For a word $w$ with $N=\binom{m}{2}$ ordered constituent pairs, Phase 2 uses
\begin{align}
\widehat h(w)&=\sum_i h_i+N^{q-1}\sum_{i<j}R_{ij},\\
\widehat a(w)&=\clip\!\left(c+(a_0-c)e^{-\widehat h(w)},0,1\right).
\label{eq:hazard}
\end{align}
Here $q=0$ means an average pair correction and $q=1$ means a full pair sum.
The development bank selected $q=0.6$ before confirmation, so pair damage grows
sub-extensively: faster than an average but slower than a quadratic pair sum.

Phase 2 compares two regressions:
\begin{align}
g_0 &= \text{clean accuracy + log parameter count + epoch + architecture
+ recipe} \notag\\
&\quad +\text{word length + the six primitive counts},\\
g_1 &= g_0 + \text{product, symmetric-hazard, ordered-hazard, and
max-damage forecasts}.
\end{align}
The central question is whether $g_1$ adds predictive value beyond $g_0$ on
new models and sealed words.

% =============================================================================
\section{Literature review}
% =============================================================================

\subsection{Robustness benchmarks are not performance forecasters}

ImageNet-C/CIFAR-C and ImageNet-P standardized evaluation under common
corruptions and perturbations, with metrics such as Corruption Error, mCE,
relative mCE, flip rate, and top-5 distance
\cite{hendrycks2019}.  RobustBench standardized robust model loading and
comparison \cite{croce2021}, while 3D Common Corruptions introduced more
geometry-aware failures \cite{kar2022}.  These works define important target
quantities and realistic model banks.  However, they normally measure a model
on labeled target data; they do not forecast an unavailable ordered
composition from its factors.

AugMix composes stochastic augmentation chains during training to improve
robustness and calibration \cite{hendrycks2020}.  This is close in its use of
compositions but different in objective: AugMix changes the classifier so that
it is robust, whereas the present project freezes the classifier and predicts
its future behavior.

\subsection{OOD performance prediction usually sees the target}

ATC estimates target accuracy from confidence or entropy thresholds applied to
actual unlabeled target predictions \cite{garg2022}.  Agreement-on-the-Line
uses agreement among models on actual target samples \cite{baek2022}; later
work extends this idea to foundation models and test-time adaptation
\cite{foundationagreement2024,ttaagreement2024}.  ODP-Bench places ten such
predictor families across 1,444 models and 29 OOD datasets, emphasizing
cross-model Spearman correlation \cite{yu2025}.

These are the closest works in task---predicting model performance---but they
operate in \Rthree{}.  Our intended regime is \Rtwo{}: the composed target
samples do not yet exist.  A fair paper must report the price of this stricter
information boundary, not pretend the methods receive equal inputs.

\subsection{Source accuracy is a strong confounder}

Accuracy-on-the-Line shows that in-distribution accuracy can strongly explain
OOD accuracy across model collections \cite{miller2021}.  This makes a raw
across-model correlation insufficient.  Our model-bank experiments therefore
need $g_0$, controlled rank after removing source accuracy/family/recipe, and
model-selection regret.  Recent benchmark critiques further warn that pooled
ID/OOD lines can hide split-level failure \cite{dgmisspecified2025}.

\subsection{The closest ordered-composition protocol}

Mason et al. study compositional robustness under six ordered primitives:
Contrast, Gaussian Blur, Impulse Noise, Invert, Rotate90, and Swirl
\cite{mason2023}.  Their protocol contains identity, six singles, all 30
ordered distinct pairs, and sampled words through length six---167 conditions
per dataset.  Their goal is to train and compare robust classifiers, not to
forecast a frozen classifier's behavior without target compositions.  We reuse
their independently designed word list as a falsification bridge.  This bridge
was scientifically important because it rejected our original primary method
under both fixed and severity-calibrated settings.

\subsection{Natural shifts and compositional semantics}

WILDS, DomainBed, BREEDS, and MetaShift provide natural domain, group, context,
and subpopulation shifts \cite{koh2021,gulrajani2021,santurkar2021,liang2022}.
These benchmarks are essential for external validity, but many metadata
intersections are commutative.  A region-by-year group should not be described
as a physical non-commutative operator.  A defensible extension separates
commutative factor intersections from physically ordered natural-plus-operator
words.

Sparse Differentiable Tree Machines study compositional generalization under
lexical, structural, length, template, and maximum-compound-divergence splits
\cite{soulos2024}.  This provides broad compositional-generalization context,
but it is a symbolic sequence/tree task rather than image-corruption
performance forecasting.

\subsection{Invariant learning is adjacent, not the same task}

Invariant Risk Minimization and related domain-generalization methods attempt
to learn predictors that remain stable across environments
\cite{arjovsky2019}.  They modify training and representations.  Our task is
complementary: take a frozen model, including one that may not be invariant,
and estimate its behavior under an unavailable composition.  A forecasting
paper should not claim that invariance methods ``fail'' in general; it should
show that they answer a different question.

% =============================================================================
\section{Paper-by-paper comparison and required proof of difference}
% =============================================================================

\begin{landscape}
\scriptsize
\setlength{\tabcolsep}{3pt}
\renewcommand{\arraystretch}{1.16}
\begin{longtable}{P{3.1cm}P{3.4cm}P{4.4cm}P{4.4cm}P{5.0cm}}
\caption{Closest literature, overlap, distinction, and evidence required.}
\label{tab:literature}\\
\toprule
\textbf{Paper / year} & \textbf{What it studies} & \textbf{What is similar to our work} & \textbf{How our task differs} & \textbf{What we must prove to justify the difference} \\
\midrule
\endfirsthead
\toprule
\textbf{Paper / year} & \textbf{What it studies} & \textbf{What is similar to our work} & \textbf{How our task differs} & \textbf{Required evidence} \\
\midrule
\endhead
\midrule
\multicolumn{5}{r}{\emph{Continued on next page}}\\
\endfoot
\bottomrule
\endlastfoot

Hendrycks \& Dietterich, \emph{Common Corruptions and Perturbations}, 2019 \cite{hendrycks2019}
& Standardized corruption/perturbation robustness evaluation.
& Same broad object: model behavior under image corruption; supplies mCE, relative mCE, flip-rate metrics.
& Measures labeled targets, mainly individual corruption/severity conditions; does not forecast unavailable ordered words.
& Predict standard accuracy/error and eventually mCE contributions on held-out compositions; use common corruption definitions and disclose deviations. \\

Hendrycks et al., \emph{AugMix}, 2020 \cite{hendrycks2020}
& Trains robust, calibrated classifiers with mixed augmentation chains.
& Uses composed image transformations and recognizes that mixtures improve robustness training.
& Changes classifier training; ours freezes the classifier and forecasts failure before the composed target is available.
& Evaluate frozen standard and robust models; show value is predictive rather than an artifact of retraining; include AugMix/robust recipes in the bank. \\

Arjovsky et al., \emph{Invariant Risk Minimization}, 2019 \cite{arjovsky2019}
& Learns invariant predictors across environments.
& Motivated by deployment under distribution change.
& Attempts to remove environment-specific dependence; does not represent and forecast ordered operators for a frozen model.
& Demonstrate prediction for models with different training recipes, control for source accuracy, and avoid claiming that forecasting replaces robust training. \\

Miller et al., \emph{Accuracy on the Line}, 2021 \cite{miller2021}
& Relates ID and OOD accuracy across model collections.
& Directly relevant source-only performance baseline and across-model unit.
& Does not use primitive/pair composition information or ordered words.
& Show $g_1$ improves over clean accuracy, family, recipe, epoch, and parameter count; report controlled rank and selection regret with clustered intervals. \\

Koh et al., \emph{WILDS}, 2021 \cite{koh2021}
& Natural in-the-wild distribution shifts and worst-group evaluation.
& Supplies realistic domains and deployment metrics.
& Natural groups are not generally sequential image operators; benchmark evaluates rather than forecasts compositions.
& Test commutative metadata intersections separately from natural-plus-physical ordered shifts on at least two datasets. \\

Gulrajani \& Lopez-Paz, \emph{DomainBed}, 2021 \cite{gulrajani2021}
& Standardized domain-generalization algorithms and model selection.
& Emphasizes controlled protocols and held-out domains.
& Trains domain-generalized models rather than predicting behavior of arbitrary frozen models.
& Use architecture/recipe holdouts and fixed selection protocols; show composition score transfers beyond seen model families. \\

Santurkar et al., \emph{BREEDS}, 2021 \cite{santurkar2021}
& Subpopulation shift benchmarks.
& Relevant natural factor and model-bank test source.
& Does not define physical ordered composition or target-free behavior forecasting.
& Forecast subpopulation intersections without inventing order; demonstrate calibrated error on natural factors. \\

Croce et al., \emph{RobustBench}, 2021 \cite{croce2021}
& Standardized robust model zoo and evaluation.
& Provides architecture/training diversity needed for model ranking.
& Primarily evaluates adversarial/common-corruption robustness with target labels.
& Use public preprocessing-correct checkpoints, hashes, and licenses; demonstrate incremental value across standard, corruption-robust, and adversarial recipes. \\

Garg et al., \emph{Leveraging Unlabeled Data / ATC}, 2022 \cite{garg2022}
& Estimates OOD accuracy from actual unlabeled target confidence.
& Same output: target accuracy estimate.
& ATC sees target images (\Rthree{}); our primary forecaster does not (\Rtwo{}).
& Enforce code-level no-target access; compare equal regimes; report the error cost of not seeing target images using official ATC. \\

Baek et al., \emph{Agreement-on-the-Line}, 2022 \cite{baek2022}
& Predicts OOD performance from agreement among models on target samples.
& Motivates cross-model ranking and selection.
& Requires target samples and a model ensemble; no primitive/pair factorization.
& Build a diverse model bank; report per-environment model rank and regret; compare with Agreement only as target-available \Rthree{}. \\

Kar et al., \emph{3D Common Corruptions}, 2022 \cite{kar2022}
& Geometry-aware realistic corruption benchmark.
& Supplies physically richer primitive shifts and five severities.
& Evaluates observed corruptions rather than unseen words.
& Show transfer beyond low-level synthetic CIFAR operations, with deterministic generation and standard metrics. \\

Liang \& Zou, \emph{MetaShift}, 2022 \cite{liang2022}
& Contextual distribution shifts and conflicts over many natural subsets.
& Provides factor annotations for compositional groups.
& Factors are semantic/contextual and usually commutative, not chronological operators.
& Define a mathematically appropriate commutative track and a separate natural-plus-operator ordered track. \\

Mason et al., \emph{Modularity Trumps Invariance}, 2023 \cite{mason2023}
& Trains/analyzes models under ordered corruption compositions through length six.
& Closest exact word protocol: same six primitives, 30 pairs, 167 conditions.
& Their endpoint is robustness/invariance after training; ours forecasts a frozen model using primitive/pair calibration.
& Beat strong max, retention, pair-residual, and sequence baselines on an independent word protocol. Our Phase 1 and 1B currently fail this method requirement and must remain disclosed. \\

CorruptionBenchCV, 2023 \cite{corruptionbench2023}
& Unifies C/C-bar/P/3DCC evaluation and robustness/calibration metrics.
& Defines the quantities a practical forecast should recover.
& Measures observed target data; no pre-deployment factor composition.
& Forecast per-severity accuracy, ECE, relative robustness, and model ordering under compatible transformations. \\

\emph{Foundation Models via Agreement-on-the-Line}, 2024 \cite{foundationagreement2024}
& Agreement-based prediction for lightly fine-tuned foundation models.
& Highlights the role of heterogeneous model banks and ensemble diversity.
& Uses target samples and agreement, not primitive profiles or ordered composition.
& Demonstrate value with heterogeneous backbones and compare against target-available agreement under identical model pools. \\

\emph{Test-Time Adaptation Induces Stronger Agreement-on-the-Line}, 2024 \cite{ttaagreement2024}
& Uses adaptation to improve agreement/performance estimation.
& Shares deployment selection motivation.
& Changes models at target time and consumes target images; ours freezes models and forecasts before target arrival.
& Keep models frozen in the primary claim; treat adaptation as a separate operational regime. \\

Soulos et al., \emph{Sparse Tree Operations}, 2024 \cite{soulos2024}
& Neurosymbolic sequence/tree compositional generalization.
& Tests length and structural extrapolation.
& Symbolic grammar and parse trees are not image-shift operators or performance estimates.
& Include matched sequence baselines (GRU/Transformer/MLP) and show gains are not generic sequence-model capacity. \\

Salaudeen et al., \emph{DG Benchmarks with Accuracy on the Line}, 2025 \cite{dgmisspecified2025}
& Audits whether benchmark conclusions are dominated by ID/OOD lines.
& Direct warning for our Phase-2 model-bank claim.
& Audits benchmark validity rather than proposing composition forecasting.
& Report split-level results, source-controlled rank, family/recipe holdouts, and avoid one pooled correlation as proof. \\

Yu et al., \emph{ODP-Bench}, 2025 \cite{yu2025}
& Ten OOD predictors over 1,444 models and 29 datasets.
& Closest umbrella benchmark for performance prediction; uses model-rank Spearman.
& Existing predictors see actual unlabeled targets; benchmark has no ordered pre-deployment composition track.
& Scale beyond the CPU pilot, use official algorithms, reproduce ODP-style ranking, and establish the new information regime without claiming direct SOTA prematurely. \\

\end{longtable}
\end{landscape}

\subsection{Provisional novelty statement}

The defensible working statement is:
\begin{quote}
Existing OOD performance prediction estimates performance after unlabeled target
samples become available, while compositional-robustness work measures or
improves classifiers on observed compositions.  We study the earlier
pre-deployment problem of forecasting an unavailable ordered target composition
from characterized primitives and limited pair calibration.
\end{quote}
This is provisional.  Before submission, forward/backward citation searches
must be repeated in Semantic Scholar, OpenAlex, Google Scholar, OpenReview, and
conference proceedings.  The project must not claim to have invented
semigroups, matrix multiplication, commutators, $L_p$ norms, hazards, or
mixture-of-experts.

% =============================================================================
\section{What must be proved empirically and mathematically}
% =============================================================================

\subsection{Empirical obligations}

A credible difference from prior work requires all of the following:
\begin{enumerate}
  \item \textbf{Target-access proof.}  Code and artifacts must show that target
  composition images/labels do not enter \Rtwo{} fitting or prediction.
  \item \textbf{Incremental value.}  Composition-aware $g_1$ must beat $g_0$
  containing clean accuracy, architecture, training recipe, parameter count,
  epoch, length, and primitive identity, with a family/lineage-clustered
  confidence interval excluding zero.
  \item \textbf{Strong same-regime baselines.}  Compare max damage, additive,
  product retention, signed $L_p$, pair ANOVA/residual, and matched neural
  sequence models using the same primitive/pair information.
  \item \textbf{Across-model utility.}  Report model-rank Spearman, controlled
  rank, top-$k$ selection regret, and near-oracle selection rate.
  \item \textbf{Breadth.}  Test held-out model families, recipes, severities,
  datasets, and realistic corruptions; one weak CIFAR bank is not sufficient.
  \item \textbf{R2 versus R3 clarity.}  Compare official target-image methods
  separately and report the price of not seeing the target.
  \item \textbf{Failure disclosure.}  Retain Phase-1/1B negative results and
  report clipping, saturation, support distance, and per-length behavior.
  \item \textbf{Order claim discipline.}  Claim non-commutativity only if the
  ordered model beats a symmetric capacity-matched model on preregistered
  reverse pairs across datasets and model families.  Current Phase-2 evidence
  does not satisfy this obligation.
\end{enumerate}

\subsection{Mathematical obligations}

The implementation-matched theory should establish:
\begin{enumerate}
  \item exact inverse consistency of the numerical $\phi_p,\phi_p^{-1}$ pair;
  \item associativity of the accumulated state, not automatically of every
  nonlinear decoder heuristic;
  \item permutation invariance of $\Gamma$ and reversal antisymmetry of
  $\Omega$;
  \item equivalence of matrix and closed-form extraction;
  \item a length-dependent error decomposition of the form
  \[
  \|\widehat B(w)-B(w)\|
  \le C_1m\epsilon_{\mathrm{primitive}}
  +C_2\binom{m}{2}\epsilon_{\mathrm{pair}}
  +C_3\epsilon_{\mathrm{higher-order}};
  \]
  \item an identifiability limitation: primitive marginals alone cannot
  determine arbitrary pair behavior;
  \item the expressivity limit of a transitive potential order---it cannot
  represent arbitrary cyclic pair preferences.
\end{enumerate}

% =============================================================================
\section{Experimental evidence: complete and honest progression}
% =============================================================================

\subsection{V1: the original four-primitive CIFAR experiment failed}

The original CIFAR-10 study used 45 observed conditions and 33 held-out
conditions.  It is important not to replace this first result with later
improvements.

\begin{center}
\begin{tabular}{lrrrr}
\toprule
\textbf{Method} & \textbf{BPE $\downarrow$} & \textbf{Accuracy MAE $\downarrow$} & \textbf{Rank $\rho$ $\uparrow$} & \textbf{Interpretation}\\
\midrule
Cross-attention M3 & \textbf{0.236} & 7.96pp & \textbf{0.881} & Best zero-shot method\\
No-composition M0 & 0.381 & \textbf{7.62pp} & 0.821 & Strong simple baseline\\
No-order algebra M6 & 0.536 & 7.76pp & 0.647 & Better than full M5\\
Original Shift Algebra M5 & 0.962 & 11.62pp & 0.536 & Primary hypothesis failed\\
Additive behavior & 1.828 & 34.21pp & 0.708 & Severe over-composition\\
\bottomrule
\end{tabular}
\end{center}

M5 beat the additive and reconstruction-only controls, but it was significantly
worse than the no-composition, MLP, and attention baselines.  It also predicted
the flagship Blur/Noise order direction incorrectly.  This motivated smaller,
behavior-anchored, exactly associative V2 models.

\subsection{V2/V3: positive but development-cycle CIFAR evidence}

The six-primitive V2 CIFAR-10 benchmark used 155 observed conditions and 123
held-out conditions.  Later reviewer-proofing compared 14 methods and added
primitive holdouts, a second architecture, external datasets, counterfactuals,
and clustered inference.

\begin{center}
\small
\begin{tabular}{lrrrrrr}
\toprule
\textbf{Method} & \textbf{BPE} & \textbf{Triple BPE} & \textbf{Order BPE} & \textbf{Acc. MAE} & $\bm{\rho_m}$ & \textbf{Order sign}\\
\midrule
$M_B$ operator algebra & \textbf{0.0152} & \textbf{0.0122} & 0.0297 & 4.24pp & \textbf{0.929} & \textbf{0.90}\\
$M_A$ behavior algebra & 0.0164 & 0.0152 & \textbf{0.0234} & \textbf{4.23pp} & 0.915 & \textbf{0.90}\\
Power+offset & 0.0285 & 0.0260 & 0.0331 & 6.44pp & 0.924 & 0.00\\
Signed $L_p$ only & 0.0289 & 0.0262 & 0.0351 & 6.46pp & 0.926 & 0.00\\
Low-rank bilinear & 0.0301 & 0.0215 & 0.0662 & 4.75pp & 0.922 & 0.85\\
Matched MLP & 0.0344 & 0.0234 & 0.0761 & 5.74pp & 0.872 & 0.82\\
Mean pooling & 0.0544 & 0.0530 & 0.0545 & 8.31pp & 0.925 & 0.00\\
Additive behavior & 0.5674 & 0.6040 & 0.3826 & 11.48pp & 0.919 & 0.00\\
\bottomrule
\end{tabular}
\end{center}

Ablations were supportive: removing the $L_p$ law increased BPE from 0.0152 to
0.1170, and removing the commutator increased order BPE from 0.0297 to 0.0912
and reduced sign match from 0.90 to 0.16.  However, these V1--V3 experiments
were part of an iterative development cycle, used custom BPE, and mainly
ranked environments for one model.  They are useful development evidence, not
an untouched final confirmation or an ODP-Bench-equivalent result.

\subsection{Phase 1: independent fixed-severity bridge failed}

Phase 1 reused the independently designed Mason word list with 9,000 CIFAR-10
and 16,920 EMNIST confirmation images.  All six singles and 30 ordered pairs
were calibration inputs; 130 words of lengths three through six were targets.

\begin{center}
\begin{tabular}{lrrr}
\toprule
\textbf{Method} & \textbf{CIFAR-10 MAE} & \textbf{EMNIST MAE} & \textbf{Dataset macro}\\
\midrule
Max-absolute primitive envelope & \textbf{4.664pp} & \textbf{0.457pp} & \textbf{2.561pp}\\
Pair residual & 6.652pp & 2.262pp & 4.457pp\\
Potential Shift Algebra & 7.081pp & 2.267pp & 4.674pp\\
Signed $L_p$ & 7.214pp & 2.270pp & 4.742pp\\
\bottomrule
\end{tabular}
\end{center}

Potential Shift Algebra was 82.5\% worse than the max envelope.  The reason was
severe floor saturation: mean CIFAR accuracy fell from 23.69\% at length one to
7.76\% at length six, while EMNIST was near its 2.13\% chance level from short
words onward.  This is a transparent failed method gate.

\subsection{Phase 1B: severity calibration changed the winner, not the verdict}

Phase 1B attenuated each primitive to approximately an 8\% normalized
clean-to-chance single-shift drop and froze strengths before target evaluation.
The max envelope ceased to be competitive, confirming that its Phase-1 victory
was severity dependent.  Nevertheless, all methods underestimated accumulated
higher-order damage.

\begin{center}
\begin{tabular}{lrrr}
\toprule
\textbf{Method} & \textbf{CIFAR-10 MAE} & \textbf{EMNIST MAE} & \textbf{Dataset macro}\\
\midrule
Pair residual & \textbf{17.202pp} & \textbf{20.473pp} & \textbf{18.837pp}\\
Pair-feature ridge & 17.295pp & 20.475pp & 18.885pp\\
Potential Shift Algebra & 17.765pp & 23.385pp & 20.575pp\\
Product/additive-style fixed laws & \multicolumn{3}{c}{substantially worse; see frozen report}\\
\bottomrule
\end{tabular}
\end{center}

Potential algebra was 9.22\% worse than pair residual in macro MAE.  Its order
channel did improve a small reversal diagnostic, but the dynamic-range and
primary method gates failed.  The mechanistic lesson was that an
interaction average is too \emph{intensive}: it remains bounded as word length
grows even when total damage accumulates.

% =============================================================================
\section{Phase 2: what was done, in simple mathematical terms}
% =============================================================================

\subsection{Why Phase 2 changed the evaluation unit}

Earlier phases mostly asked:
\begin{quote}
For one classifier, can we rank and predict many composed environments?
\end{quote}
Phase 2 asks the OOD-prediction-style question:
\begin{quote}
For one composed environment, can we rank many classifiers and select a good
model before running that environment?
\end{quote}
These are different estimands.  Phase 2 created:
\begin{itemize}
  \item 30 development snapshots from 15 independent training lineages;
  \item 30 confirmation snapshots from 15 new lineages;
  \item five lightweight architecture families and three training recipes;
  \item 37 clean/single/pair calibration conditions;
  \item 40 method-development words and 90 sealed confirmation words;
  \item disjoint 500-image development and 1,000-image confirmation behavior
  partitions.
\end{itemize}
Epoch-one and epoch-two snapshots share a lineage, so inference clusters them
rather than pretending there are 30 independent training runs.

\subsection{Step-by-step Phase-2 prediction}

For each model $f$:
\begin{enumerate}
  \item Measure clean accuracy $a_0(f)$.
  \item Measure six primitive accuracies $a_i(f)$.
  \item Measure all 30 ordered pair accuracies $a_{i\to j}(f)$.
  \item Convert primitive and pair accuracies to hazards.
  \item Compute ordered pair residuals $R_{ij}$.
  \item Compose a new word using Equation~\eqref{eq:hazard} with frozen
  $q=0.6$.
  \item Give these forecasts, together with source/model/word metadata, to
  $g_1$.
  \item Compare against $g_0$, which receives the same metadata but no
  composition forecast.
\end{enumerate}
The exponent, ridge coefficients, feature scales, metrics, and gates were
committed before confirmation lineages were trained/evaluated.

\subsection{A small numerical example}

Suppose clean accuracy is $a_0=0.80$ and chance is $c=0.10$.  Blur accuracy is
$0.70$ and Noise accuracy is $0.65$.  Their retentions and hazards are
\begin{align}
r_B&=\frac{0.70-0.10}{0.80-0.10}=0.857,
&h_B&=-\log(0.857)=0.154,\\
r_N&=\frac{0.65-0.10}{0.80-0.10}=0.786,
&h_N&=-\log(0.786)=0.241.
\end{align}
If measured Blur$\to$Noise accuracy is 0.55, then
\[
h_{B\to N}=-\log\left(\frac{0.55-0.10}{0.80-0.10}\right)=0.442
\]
and
\[
R_{BN}=0.442-0.154-0.241=0.047.
\]
A longer word sums primitive hazards and scaled pair residuals, then converts
hazard back to bounded accuracy.  Unlike raw additive accuracy loss, the output
is naturally tied to clean and chance performance.

\subsection{Phase-2 sealed results}

On 30 confirmation snapshots and 90 sealed words, all preregistered incremental
criteria passed.

\begin{center}
\begin{tabular}{lrr}
\toprule
\textbf{Metric} & $\bm{g_0}$ \textbf{source/model/word} & $\bm{g_1}$ \textbf{+ composition}\\
\midrule
Accuracy MAE & 3.304pp & \textbf{1.737pp}\\
Raw model-rank Spearman & 0.242 & \textbf{0.665}\\
Source/family/recipe-controlled Spearman & $-0.235$ & \textbf{0.657}\\
Top-1 model-selection regret & 2.158pp & \textbf{0.489pp}\\
Top-3 model-selection regret & 1.907pp & \textbf{0.214pp}\\
Selected model within 1pp of oracle & 44.4\% & \textbf{78.9\%}\\
\bottomrule
\end{tabular}
\end{center}

The direct MAE reduction is 47.4\%.  A 10,000-draw paired bootstrap clustered
by reverse-word group and training lineage estimated a 1.529pp improvement,
with 95\% interval $[0.981,2.191]$pp.

\begin{center}
\begin{tabular}{lrrr}
\toprule
\textbf{Forecast} & \textbf{Information} & \textbf{MAE} & \textbf{Model-rank $\rho$}\\
\midrule
$g_1$ metadata + composition & \Rtwo{} & \textbf{1.737pp} & \textbf{0.665}\\
Symmetric hazard ($q=0.6$) & \Rtwo{} & 2.089pp & 0.610\\
Ordered hazard ($q=0.6$) & \Rtwo{} & 2.093pp & 0.607\\
$g_0$ source/model/word & \Rzero{} + word & 3.304pp & 0.242\\
Product retention & \Rone{} & 5.226pp & 0.249\\
Accuracy-space pair residual & \Rtwo{} & 7.265pp & 0.526\\
Max primitive damage & \Rone{} & 12.033pp & 0.119\\
ATC (simplified pilot implementation) & \Rthree{} & 20.495pp & $-0.180$\\
DoC (simplified pilot implementation) & \Rthree{} & 19.247pp & $-0.067$\\
\bottomrule
\end{tabular}
\end{center}

The ATC/DoC numbers must not be generalized to official ODP-Bench: the pilot
used weak, poorly calibrated two-epoch models and simplified declared
implementations.  The accuracy-space pair residual also achieved slightly
lower top-1 regret than $g_1$ despite worse MAE, showing that no method dominates
every endpoint.

\subsection{What Phase 2 did not prove}

The symmetric hazard model slightly beat the ordered model:
\[
\MAE_{\mathrm{sym}}=2.089\text{pp}
<\MAE_{\mathrm{ord}}=2.093\text{pp}.
\]
On the reversal diagnostic, symmetric MAE was 0.570pp versus ordered 0.582pp,
and ordered sign accuracy was 65\%.  Therefore Phase 2 supports extensive
composition-aware forecasting, not a general non-commutative order claim.

The bank also contains only 15 independent confirmation lineages, five
lightweight architecture families, three recipes, CIFAR-10 only, and models
with 29.8--53.5\% clean accuracy.  It is a positive CPU pilot, not the planned
75+ model final benchmark.

\subsection{Phase 2B: adaptive mixture on a fresh CPU confirmation}

After the Phase-2 result, the project tested an Adaptive Composition Hazard
Mixture (ACHM) that forms a convex combination of dominant-shift, independent,
symmetric, and ordered hazard experts.  The gate was fitted only on the earlier
development roles and frozen before training 15 new lineages.  Confirmation
used 30 new snapshots, a disjoint 1,000-image behavior partition, and 64 words
that did not occur in the earlier 167-word protocol.

ACHM reduced MAE from 1.719pp for the preregistered ordered-hazard comparator
to 1.593pp, a 7.3\% relative reduction.  However, the clustered 95\% interval
for the 0.126pp improvement was $[-0.025,0.295]$pp, and controlled model-rank
Spearman declined from 0.822 to 0.811.  The complete gate therefore failed.
The earlier frozen metadata-plus-composition regression remained stronger at
1.241pp MAE.  A pair-budget curve improved from 5.357pp with no pairs to
3.431pp, 1.726pp, and 1.593pp with 5, 10, and 15 bidirectional pair groups,
respectively, but no random-selection comparator was preregistered.

The adaptive gate placed 80.8\% average weight on symmetric compounding.  Its
order sign accuracy was only 61.2\%, and expert disagreement failed as an
uncertainty score.  Thus Phase 2B supplies a useful negative gate and
calibration-cost result, not evidence that ACHM supersedes the Phase-2 model.

\subsection{Phase 2C: active pair calibration under a fixed budget}

A separate fresh experiment tested the practical question of which five
bidirectional pairs should be measured.  The active design was frozen from
development interaction magnitude, word frequency, and primitive coverage.
It was compared at equal cost with twenty uniform-random five-group designs on
15 new lineages, a new 1,000-image partition, and another 64 new words.

The selected design obtained 3.176pp MAE versus 3.636pp expected random MAE,
a 0.459pp improvement with clustered 95\% interval $[0.119,0.889]$pp.
Controlled rank was 0.474 versus a random-design median of 0.417, and top-1
regret was 3.870pp versus 3.984pp.  All preregistered Phase-2C criteria passed.
The selected design beat 17 of 20 random subsets, but the best random subset
reached 2.592pp.  Full fifteen-group ACHM and prior $g_1$ reached 1.603pp and
1.410pp, respectively, using three times as many directed pair measurements.
The result therefore supports selection over expected random calibration, not
global optimality or full-budget equivalence.

% =============================================================================
\section{Current scientific conclusion and recommended framing}
% =============================================================================

The combined evidence does not support a universal Shift Algebra.  Instead, it
shows multiple composition regimes:
\begin{enumerate}
  \item Under severe floor saturation, a dominant-shift envelope can be best.
  \item Under calibrated but rapidly accumulating damage, pair-average
  interaction laws can be too optimistic.
  \item Across a model bank, sub-extensive hazard composition plus model/word
  metadata can add substantial predictive value.
  \item Order features can work on selected mechanisms but are not yet broadly
  reliable.
  \item Under a constrained calibration budget, development-informed pair
  selection can outperform expected random pair selection, although it remains
  materially weaker than complete pair calibration.
\end{enumerate}

The safest paper-level framing is therefore:
\begin{quote}
\textbf{Pre-deployment compositional performance forecasting:} characterize
when max, independent-retention, and pair-compounding laws are useful, and test
whether their forecasts add value beyond model and word metadata without
observing the composed target.
\end{quote}
Shift Algebra can remain an interpretable order-aware expert, but the current
headline should be composition-aware forecasting, not a claim that all shifts
obey one non-commutative algebra or that adaptive law selection is already
established.

Phase 2B tested the natural convex mixture
\[
\widehat h(w)=
\alpha_{\max}h_{\max}(w)+
\alpha_{\mathrm{prod}}h_{\mathrm{prod}}(w)+
\alpha_{\mathrm{sym}}h_{\mathrm{sym}}(w)+
\alpha_{\mathrm{ord}}h_{\mathrm{ord}}(w),
\quad
\alpha_k\ge0,\quad\sum_k\alpha_k=1,
\]
using only pre-target gate inputs.  The small MAE gain did not satisfy the full
gate.  Any revised adaptive model must treat all Phase-2B outcomes as opened
development evidence and use another genuinely fresh test; it must not be
rescued by tuning on this confirmation.

% =============================================================================
\section{Checklist before a NeurIPS submission}
% =============================================================================

\begin{enumerate}
  \item Complete a refreshed multi-index novelty search and update the claim
  matrix; do not claim ``first'' until overlap is audited.
  \item Do not tune ACHM on Phase-2B confirmation; freeze any revision on a
  separate development set before another fresh confirmation.
  \item Evaluate at least 75 final checkpoints, six architecture families,
  four recipe families, and CIFAR-10/CIFAR-100, preferably using public models.
  \item Add architecture-family and recipe holdouts rather than seed-only
  transfer.
  \item Integrate official ODP predictors and clearly label \Rtwo{} versus
  \Rthree{} access.
  \item Compare active pair selection with random, magnitude-only, coverage-only,
  and D-optimal designs on a new held-out architecture/dataset.
  \item Add uncertainty, support detection, and abstention with coverage tests.
  \item Test realistic corruptions and at least two natural/hybrid settings.
  \item Report all negative phases and all primary endpoints, not only the best
  method/metric.
  \item Obtain an independent mathematical and leakage audit.
\end{enumerate}

% =============================================================================
\section{Reproducibility map}
% =============================================================================

\begin{center}
\begin{tabular}{P{5.0cm}P{9.5cm}}
\toprule
\textbf{Artifact} & \textbf{Purpose}\\
\midrule
\texttt{MATHEMATICAL\_TUTORIAL.tex} & Full implementation-faithful derivation of historical and V4 algebra.\\
\texttt{NEURIPS\_DEVELOPMENT\_PLAN.md} & Benchmark, literature, metric, compute, and gate roadmap.\\
\texttt{RESULTS.md} & Original V1 negative CIFAR result.\\
\texttt{RESULTS\_V2.md} and \texttt{REVIEWER\_PROOFING\_V3.md} & Historical development-cycle positive CIFAR/external evidence.\\
\texttt{v4/phase1/PHASE1\_FINAL\_REPORT.md} & Independent fixed-severity long-word failure.\\
\texttt{v4/phase1b/PHASE1B\_FINAL\_REPORT.md} & Severity-calibrated long-word failure.\\
\texttt{v4/phase2/PHASE2\_PILOT\_FINAL\_REPORT.md} & Sealed positive CIFAR model-bank pilot.\\
\texttt{v4/phase2/INFORMATION\_AUDIT.md} & Post-confirmation leak and correctness audit.\\
\texttt{v4/phase2b/PHASE2B\_FINAL\_REPORT.md} & Fresh adaptive-mixture and pair-budget results, including the failed primary gate.\\
\texttt{v4/phase2b/INFORMATION\_AUDIT.md} & Phase-2B chronology, information-access, clustering, and numerical audit.\\
\texttt{v4/phase2c/PHASE2C\_FINAL\_REPORT.md} & Fresh active-versus-random pair-calibration results.\\
\texttt{v4/phase2c/INFORMATION\_AUDIT.md} & Phase-2C access, equal-budget, clustering, and numerical audit.\\
\bottomrule
\end{tabular}
\end{center}

% =============================================================================
\begin{thebibliography}{99}
% =============================================================================

\bibitem{arjovsky2019}
M. Arjovsky, L. Bottou, I. Gulrajani, and D. Lopez-Paz.
\newblock Invariant Risk Minimization.
\newblock arXiv preprint, 2019.

\bibitem{hendrycks2019}
D. Hendrycks and T. Dietterich.
\newblock Benchmarking Neural Network Robustness to Common Corruptions and Perturbations.
\newblock ICLR, 2019.
\newblock \url{https://github.com/hendrycks/robustness}.

\bibitem{hendrycks2020}
D. Hendrycks et al.
\newblock AugMix: A Simple Data Processing Method to Improve Robustness and Uncertainty.
\newblock ICLR, 2020.
\newblock \url{https://github.com/google-research/augmix}.

\bibitem{miller2021}
J. Miller et al.
\newblock Accuracy on the Line: On the Strong Correlation Between Out-of-Distribution and In-Distribution Generalization.
\newblock ICML, 2021.
\newblock \url{https://github.com/millerjohnp/linearfits_app}.

\bibitem{koh2021}
P. W. Koh et al.
\newblock WILDS: A Benchmark of in-the-Wild Distribution Shifts.
\newblock ICML, 2021.
\newblock \url{https://github.com/p-lambda/wilds}.

\bibitem{gulrajani2021}
I. Gulrajani and D. Lopez-Paz.
\newblock In Search of Lost Domain Generalization.
\newblock ICLR, 2021.
\newblock \url{https://github.com/facebookresearch/DomainBed}.

\bibitem{santurkar2021}
S. Santurkar et al.
\newblock BREEDS: Benchmarks for Subpopulation Shift.
\newblock ICLR, 2021.
\newblock \url{https://github.com/MadryLab/BREEDS-Benchmarks}.

\bibitem{croce2021}
F. Croce et al.
\newblock RobustBench: A Standardized Adversarial Robustness Benchmark.
\newblock NeurIPS Datasets and Benchmarks, 2021.
\newblock \url{https://github.com/RobustBench/robustbench}.

\bibitem{garg2022}
S. Garg et al.
\newblock Leveraging Unlabeled Data to Predict Out-of-Distribution Performance.
\newblock ICLR, 2022.
\newblock \url{https://github.com/saurabhgarg1996/ATC_code}.

\bibitem{baek2022}
C. Baek et al.
\newblock Agreement-on-the-Line: Predicting the Performance of Neural Networks under Distribution Shift.
\newblock NeurIPS, 2022.
\newblock \url{https://github.com/kebaek/Agreement-on-the-line}.

\bibitem{kar2022}
O. F. Kar et al.
\newblock 3D Common Corruptions and Data Augmentation.
\newblock CVPR Oral, 2022.
\newblock \url{https://github.com/EPFL-VILAB/3DCommonCorruptions}.

\bibitem{liang2022}
W. Liang and J. Zou.
\newblock MetaShift: A Dataset of Datasets for Evaluating Contextual Distribution Shifts and Training Conflicts.
\newblock ICLR, 2022.
\newblock \url{https://github.com/Weixin-Liang/MetaShift}.

\bibitem{mason2023}
I. Mason et al.
\newblock Modularity Trumps Invariance for Compositional Robustness.
\newblock Preprint, 2023.
\newblock \url{https://github.com/ianxmason/compositional-robustness}.

\bibitem{corruptionbench2023}
CorruptionBenchCV project.
\newblock A Survey on the Robustness of Computer Vision Models against Common Corruptions / CorruptionBenchCV.
\newblock Survey and benchmark tool, 2023.
\newblock \url{https://github.com/nis-research/CorruptionBenchCV}.

\bibitem{foundationagreement2024}
\newblock Predicting the Performance of Foundation Models via Agreement-on-the-Line.
\newblock NeurIPS, 2024.
\newblock \url{https://openreview.net/forum?id=aJx9onwsR4}.

\bibitem{ttaagreement2024}
\newblock Test-Time Adaptation Induces Stronger Accuracy and Agreement-on-the-Line.
\newblock NeurIPS, 2024.
\newblock \url{https://openreview.net/forum?id=giXUx4VH9t}.

\bibitem{soulos2024}
P. Soulos et al.
\newblock Compositional Generalization Across Distributional Shifts with Sparse Tree Operations.
\newblock NeurIPS Spotlight, 2024.
\newblock \url{https://github.com/psoulos/sdtm}.

\bibitem{dgmisspecified2025}
O. Salaudeen et al.
\newblock Are Domain Generalization Benchmarks with Accuracy on the Line Misspecified?
\newblock Preprint, 2025.
\newblock \url{https://github.com/olawalesalaudeen/misspecified_DG_benchmarks}.

\bibitem{yu2025}
H. Yu et al.
\newblock ODP-Bench: Benchmarking Out-of-Distribution Performance Prediction.
\newblock ICCV, 2025.
\newblock \url{https://github.com/h-yu16/Performance_Prediction}.

\end{thebibliography}

\end{document}
